% TOP_PROBLEM: {"id": 30006722, "title": "Grothendieck–Teichmüller Conjecture", "category_id": 1, "category": {"id": 1, "name": "number_theory", "display_name": "Number Theory", "description": "Properties of integers, prime numbers, Diophantine equations.", "slug": "number-theory", "order_index": 1, "created_at": "2026-07-31T15:26:25.670Z"}, "status": "open"}
% FIRST_POSTED: 2026-09-27
% !TeX program = lualatex
% ATTEMPT_STATUS: unresolved
\documentclass[11pt]{article}
\usepackage[margin=25mm]{geometry}
\usepackage{fontspec}
\setmainfont{Latin Modern Roman}
\usepackage{amsmath,amssymb,mathtools}
\usepackage{unicode-math}
\setmathfont{Latin Modern Math}
\usepackage{xurl,hyperref}
\hypersetup{colorlinks=true,urlcolor=blue}
\setcounter{secnumdepth}{0}
\setlength{\parindent}{0pt}
\setlength{\parskip}{5pt}
\setlength{\emergencystretch}{3em}
\begin{document}
\section*{\texorpdfstring{Grothendieck–Teichmüller Conjecture}{Grothendieck–Teichmüller Conjecture}}\label{30006722}
\subsection{Definitions and mathematical statement}
Let $\widehat F_2$ be the free profinite group on $x,y$, and let $\chi:G_{{\mathbb{Q}}}\to\widehat{{\mathbb{Z}}}^\times$ be the cyclotomic character. Drinfeld's original $\widehat{\mathrm{GT}}$ consists of the compatible invertible profinite pairs satisfying the commutator, symmetry, hexagon, and pentagon conditions; all these conditions, not only the pair equations, are incorporated from the source. Under the fixed Ihara splitting,
\[
 \mathrm{Ih}(g)=\bigl((\chi(g)-1)/2,f_g\bigr),\qquad
 x\mapsto x^{\chi(g)},\quad y\mapsto f_g^{-1}y^{\chi(g)}f_g.
\]
The question is surjectivity of this actual embedding $G_{{\mathbb{Q}}}\hookrightarrow\widehat{\mathrm{GT}}$. Gentle, pro-$\ell$, and prounipotent groups are not substituted. Composition and coordinate conventions must agree with Section 1.3 of the cited paper.
\par\smallskip\textit{Formulation and concept references:} \href{https://arxiv.org/pdf/2405.11725v3}{[S1]}.

\subsection{Short English statement}
Is every symmetry in the original profinite Grothendieck--Teichmüller group induced by a unique automorphism of the algebraic closure of the rationals?
\subsection{Research attempt}
\paragraph{Process, scope, and first gap.}
This GPT-6 Astra Ultra attempt, dated 27 September 2026, studies the actual Ihara embedding by finite quotients and compactness. It derives a cyclotomic-kernel reduction and checks the compatibility needed for a passage to the inverse limit. These are elementary structural deductions, not a claim of a new theorem identifying the two groups.

Throughout let $H=\widehat{\mathrm{GT}}$ be Drinfeld's original profinite group, with every commutator, symmetry, hexagon, pentagon, and invertibility requirement from the statement retained. Let $I:G_{\mathbb Q}\hookrightarrow H$ be the fixed Ihara map. The first missing assertion in this approach is
\[
 \forall U\triangleleft_{\mathrm{open}}H,\qquad
       I(G_{\mathbb Q})U=H.
 \tag{91.1}
\]
The calculations establish why (91.1) suffices and which tests it includes. They do not prove it for arbitrary open normal $U$. A quotient or family of quotients selected because it is tractable cannot replace this universal quantifier.

\paragraph{Source check and group conventions.}
The coordinate choice is precisely [S1, Section 1.3, equations (1.5)--(1.7)]. Its composition formula (2.6) agrees with the direct calculation below. The paper distinguishes its gentle groupoid from the original group, explicitly in Remark 1.2; it proves surjectivity onto a particular family of nonabelian quotients and a resulting infinite quotient. Those results do not assert surjectivity onto $H$. Drinfeld's original definition is in [E1, Section 4]. Nothing from the alternative formulations or claimed results listed in [S3]--[S4] is used as a proof of the target.

Write $h=(m,f)$ and $\lambda(h)=2m+1\in\widehat{\mathbb Z}^{\times}$. Its action is
\[
 \alpha_h(x)=x^{\lambda(h)},\qquad
 \alpha_h(y)=f^{-1}y^{\lambda(h)}f.
\]
The expression $(\lambda-1)/2$ is unambiguous in $\widehat{\mathbb Z}$: a unit is odd in its $2$-adic component, and multiplication by two is injective on $\widehat{\mathbb Z}$ with image the elements even at that component. One must not interpret the coordinate $m$ as the cyclotomic character itself.

For $h_i=(m_i,f_i)$, composition $\alpha_{h_1}\circ\alpha_{h_2}$ sends $y$ to
\[
 \alpha_{h_1}(f_2)^{-1}f_1^{-1}
       y^{\lambda(h_1)\lambda(h_2)}
       f_1\alpha_{h_1}(f_2).
\]
Thus in the source convention the coordinates are
\[
 h_1h_2=
 \bigl(m_1+m_2+2m_1m_2,\ f_1\alpha_{h_1}(f_2)\bigr).
 \tag{91.2}
\]
In particular $\lambda:H\to\widehat{\mathbb Z}^{\times}$ is a continuous homomorphism and $\lambda\circ I=\chi$. The derivation checks the action order; it does not claim that arbitrary pairs with this action satisfy the pentagon or belong to $H$.

\paragraph{The compactness lemma.}
The image $C=I(G_{\mathbb Q})$ is compact, since $G_{\mathbb Q}$ is compact and $I$ is continuous. Because the profinite group $H$ is Hausdorff, $C$ is closed. We claim
\[
 C=H\quad\Longleftrightarrow\quad
 C\longrightarrow H/U\text{ is onto for every open normal }U.
 \tag{91.3}
\]
The forward implication is immediate. For the reverse implication fix $h\in H$ and form closed subsets of the Galois group
\[
 F_U(h)=\{g\in G_{\mathbb Q}:I(g)\in hU\}.
\]
Each is nonempty by the assumed quotient surjectivity. Given $U_1,\ldots,U_r$, their intersection $V=\bigcap_iU_i$ is again open normal. A point in $F_V(h)$ belongs to every $F_{U_i}(h)$. Hence these closed sets have the finite-intersection property. Compactness supplies a point $g$ belonging to all of them. Since the intersection of the open normal subgroups of a profinite group is the identity, $I(g)=h$. \hfill$\square$

Only a cofinal family of open normal subgroups is needed: every open normal $U$ must contain some subgroup in the family. If the family is not explicitly closed under intersections, cofinality still supplies a member contained in any finite intersection. This proves compatibility of the witnesses; it is not necessary to choose an arbitrary finite-level witness and hope that it extends to every later level.

Equivalently, a failure of surjectivity would already be visible in one finite quotient. If $h\notin C$, closedness gives an open coset $hU$ disjoint from $C$, and $hU$ is a missing element of $H/U$. Thus a genuine counterexample has a finite obstruction, even though locating such a quotient or proving surjectivity onto all of them may be very difficult.

\paragraph{Separating off the known cyclotomic part.}
The cyclotomic character is onto. Indeed, for every positive integer $n$, the Galois action on a primitive $n$th root of unity realizes
\[
 \operatorname{Gal}(\mathbb Q(\mu_n)/\mathbb Q)
       \cong(\mathbb Z/n\mathbb Z)^{\times}.
\]
Every automorphism of this finite Galois extension extends to $\overline{\mathbb Q}$. The same compactness argument, or the inverse-limit description of the cyclotomic extension, gives surjectivity onto $\widehat{\mathbb Z}^{\times}$. Consequently $\lambda$ is onto as well.

Put
\[
 K_{\mathrm{cyc}}=\ker\chi
  =\operatorname{Gal}(\overline{\mathbb Q}/\mathbb Q(\mu_\infty)),
 \qquad H_1=\ker\lambda.
\]
Then
\[
 I(G_{\mathbb Q})=H\quad\Longleftrightarrow\quad
 I(K_{\mathrm{cyc}})=H_1.
 \tag{91.4}
\]
If $I$ is onto and $h\in H_1$, any preimage $g$ has $\chi(g)=1$. Conversely, given $h\in H$, choose $g_0$ with $\chi(g_0)=\lambda(h)$. The element $I(g_0)^{-1}h$ lies in $H_1$. If it equals $I(k)$ with $k\in K_{\mathrm{cyc}}$, then $h=I(g_0k)$. No splitting of the cyclotomic exact sequence is assumed. In particular, matching the first coordinate of every pair leaves exactly the nonabelian kernel problem, rather than proving surjectivity.

\paragraph{Abelianization loses a real part of the image.}
Because $f$ belongs to the closed commutator subgroup of $\widehat F_2$, the induced action on
\[
 \widehat F_2^{\mathrm{ab}}\cong\widehat{\mathbb Z}^{2}
\]
is multiplication by $\lambda(h)$ on both coordinates. All elements of $H_1$ therefore act identically on this abelianization. This is a limitation of this observation map, not a claim about the abstract abelianization of $H$ itself.

There are actual nonidentity elements of $I(K_{\mathrm{cyc}})$, so the limitation is not vacuous. The splitting field of $X^3-2$ has Galois group $S_3$: the polynomial is irreducible by Eisenstein, its splitting-field action on three roots is transitive, and complex conjugation exchanges its two nonreal roots while fixing the real one. A transitive subgroup of $S_3$ containing a transposition is $S_3$.

Choose $g_1,g_2\in G_{\mathbb Q}$ whose images are $(12)$ and $(23)$. Their commutator
\[
 g=g_1g_2g_1^{-1}g_2^{-1}
\]
has a nonidentity three-cycle as image in $S_3$, so $g\ne1$. Since the cyclotomic target is abelian, $\chi(g)=1$. Injectivity of the actual Ihara map gives $I(g)\ne1$, while its action on $\widehat F_2^{\mathrm{ab}}$ is the identity. Thus even complete knowledge of these abelianized actions cannot distinguish this genuine Galois element from the identity. We have not computed its conjugating word $f_g$, and do not identify it with the free-group word $[x,y]$.

\paragraph{Two finite-quotient stress tests.}
Surjectivity of separate coordinate observations need not give joint surjectivity. In $H_0=(\mathbb Z/2\mathbb Z)^2$, the diagonal subgroup
\[
 D=\{(0,0),(1,1)\}
\]
maps onto either coordinate but is proper. The pair $(0,1)$ has no common preimage in $D$. The compactness proof avoids this failure by requiring surjectivity for the joint quotient associated to an intersection of kernels. These coordinates are an abstract countertest, not asserted coordinate homomorphisms on $\widehat{\mathrm{GT}}$.

Even an infinite coherent tower may omit an entire direction. In $H'=\widehat{\mathbb Z}\times\mathbb Z/2\mathbb Z$, the proper closed subgroup $C'=\widehat{\mathbb Z}\times\{0\}$ maps onto every quotient obtained by reducing the first factor modulo $n$. It also maps onto their full inverse limit. The second factor remains unobserved. Accordingly, applying (91.3) to a known tower of nonabelian quotients of $H$ requires a proof that its kernels are cofinal; an infinite list and compatibility alone do not establish that property.

\paragraph{Verification and outcome.}
The checks covered the order of composition in (91.2), the uniqueness of the half-coordinate, the finite-intersection property, and the use of compactness in the Galois group. A finite permutation calculation verifies the nontrivial $S_3$ commutator, and finite-group enumeration verifies the diagonal countertest. Neither computation enumerates original GT elements or checks a finite shadow's pentagon and global liftability. All elements used as genuine GT elements arise from the already-defined embedding or are assumed to lie in $H$.

\textbf{UNRESOLVED.} Surjectivity is reduced to all finite quotients, equivalently to the full cyclotomic kernel. Abelianized tests provably miss nontrivial genuine elements, and a noncofinal tower cannot settle the limit. The missing step remains (91.1), or its kernel analogue with $H_1$; no non-Galois original GT element is constructed.

\subsection{Sources}
\begin{itemize}
\item[S1]Ivan Bortnovskyi, Vasily A. Dolgushev, Borys Holikov, Vadym Pashkovskyi, Accessing non-abelian quotients of the Grothendieck-Teichmueller group via elementary tools, arXiv2405.11725v3,6August2026,35pages.
\url{https://arxiv.org/pdf/2405.11725v3}.
\textit{Formulation source. Consulted 24 September 2026: Section 1.3, equations (1.5)--(1.6).}
\item[S2]Accessing non-abelian quotients of the Grothendieck-Teichmueller group via elementary tools.
\url{https://arxiv.org/html/2405.11725}.
\textit{Further reference; not a proof endorsement.}
\item[S3]Galois actions on surfaces and a higher genus Grothendieck-Teichmüller group.
\url{https://arxiv.org/html/2606.01466v1}.
\textit{Further reference; not a proof endorsement.}
\item[S4]Proving the Grothendieck–Teichmüller Conjecture for Profinite Spaces \& The Galois Grothendieck Path Integral.
\url{https://arxiv.org/html/2503.13006}.
\textit{Further reference; not a proof endorsement.}
\item[E1]V. G. Drinfeld, \emph{On quasitriangular quasi-Hopf algebras and on a group that is closely connected with $\operatorname{Gal}(\overline{\mathbb Q}/\mathbb Q)$}, Algebra i Analiz 2 (1990), 149--181; English translation, Leningrad Mathematical Journal 2 (1991), 829--860, Section 4.
\url{https://www.lamfa.u-picardie.fr/marin/docs/drinfeld.pdf}.
\textit{Original profinite definition; pair conventions cross-checked with [S1] on 27 September 2026.}
\end{itemize}
\end{document}
